% PNP RESEARCH REPORT::RECOVERY-STRENGTHEN8::2026-10-09T00:48:00Z
\documentclass[11pt]{article}
\usepackage[margin=0.86in]{geometry}
\usepackage{amsmath,amssymb,amsthm,booktabs,hyperref,xurl}
\newtheorem{theorem}{Theorem}
\newtheorem{lemma}[theorem]{Lemma}
\newcommand{\SAT}{\operatorname{SAT}}
\title{RECOVERY-STRENGTHEN8: Alias--Portfolio Flattening\\and Boundary-Only Recursive Closure}
\author{Recovery proof-strengthening agent}
\date{9 October 2026 UTC}
\begin{document}\sloppy\maketitle

\paragraph{Status.}
\textsc{Proved/Exact}: a virtual SAT-preserving alias composes with a fixed
sequential portfolio of source-first partial compilers; failed members cost
local work but execute no children, so only the winner contributes child
volume.  Recursive hereditary coverage is required at completed ordinary-SAT
boundaries, not at acyclic source-only alias nodes.  \textsc{Correction}:
Construction8's lifted-cross family defeats its PIVOT member and pure affine
handling, but a later pure-literal member decides the stripped source; it is
not hostile to the enlarged alias portfolio.  The all-input compiler and P
versus NP remain \textsc{Open}.  No unrestricted
$t+s<2\cos(\pi/7)$ endpoint is used.

\section{Claim card and quantifiers}
Fix repaired total DENSE-CNF1 and let $N$ be complete encoded input bit
length.  Fix one deterministic random-access machine.  Time and space are
simultaneous on that machine.  Discovery, rejected members, internal aliases,
replay, child materialization, indexes, buffers, ancestor frames and every
executed child occurrence count.

For each recursive boundary type $\tau$ there is one uniform phase
$A_\tau$.  On input $X$ of type $\tau$, it first evaluates a finite acyclic
internal DAG.  Internal edges may expose same-scale read-only virtual sources,
but make no SAT call and are completely charged to the phase.  The phase then
emits ordinary-SAT children $Q_1,\ldots,Q_k$ with declared boundary types and
an exact Boolean combiner.  The required coverage quantifier is
\begin{equation}
 \forall \tau\ \forall X\in D_\tau\ \forall i\le k(X):
       Q_i\in D_{\operatorname{type}(Q_i)}.               \tag{1}
\end{equation}
It does not quantify over internal virtual sources as future recursive roots.

\section{Recovered alias--portfolio theorem}
\begin{theorem}[ALIAS--PORTFOLIO-FLATTENING1]\label{thm:alias}
Fix a constant $d$.  A source-only transducer $T$, on $|X|=N$, either
returns FAIL before semantic calls or exposes an exact virtual DENSE-CNF1
instance $F$ with
\[
 \SAT(X)=\SAT(F),\qquad |F|\le N^{\alpha+o(1)}.
\]
Recognition costs $N^{a+o(1)}$ time and $N^{b+o(1)}$ retained state; each
adaptive bit access costs $N^{g+o(1)}$ time and $N^{s+o(1)}$ scratch.
For $j=1,\ldots,d$, let $M_j$ be a source-first partial member which completes
its dry run before any child execution and returns FAIL, DECIDE, or an
immutable correct PLAN.  Its dry run/replay uses
$m^{r_j+o(1)}$ work and $m^{p_j+o(1)}$ state.  A PLAN has maximum child
$m^{h_j+o(1)}$ and executed aggregate child volume $m^{v_j+o(1)}$, counting
multiplicity.

Run the $M_j$ sequentially through the alias and execute only the first
non-FAIL result.  Its exact accepted domain is
\[
 T^{-1}\!\left(\bigcup_{j=1}^{d}D_j\right),
\]
all-fail executes zero children, and the flattened ledger is
\[
 R=\max\{a,\max_j(g+\alpha r_j)\},\qquad
 B=\max\{b,s,\max_j(\alpha p_j)\},
\]
\[
 H=\max_j(\alpha h_j),\qquad V=\max_j(\alpha v_j).
\]
\end{theorem}
\begin{proof}
Substitute the exact alias accessor for each input probe made by $M_j$.
There are at most $m^{r_j+o(1)}$ probes/steps, so the simulated work is
$N^{g+\alpha r_j+o(1)}$.  Because $d$ is fixed and storage is reused,
sequential failed trials take a maximum exponent.  Source-first certification
prevents every failed member from executing a child; hence only the winning
PLAN contributes $H,V$.  The alias identity and the winner's certificate
give correctness.  Synthesized header, canonical bytes and EOF are part of
the exact accessor, not a merely semantic view.
\end{proof}

This theorem is sharper than applying STREAMCOMP1 around the whole recursive
execution: after an explicit child is emitted, solving it does not access the
parent alias.  Conversely, a compact descriptor is insufficient without
$g,s$: a prefix-parity alias implemented by rescanning the source makes a
linear-read member quadratic, and materializing a same-size alias costs
linear state.  Identical aliases and children still execute with multiplicity;
ORBITBATCH1 address equality does not erase them.

\section{Recovered boundary-only theorem}
\begin{theorem}[BOUNDARY-CLOSURE1]\label{thm:closure}
Assume every $A_\tau$ terminates its internal DAG before its first child call,
is correct on all $X\in D_\tau$, and satisfies (1).  Suppose, uniformly over
all exposed boundary edges $\tau\to\sigma$,
\[
 T_{\rm phase}(N)\le N^{r+o(1)},\quad
 S_{\rm phase}(N)\le N^{p+o(1)},\quad
 \max_i |Q_i|\le N^{h+o(1)},\quad
 \sum_i|Q_i|\le N^{v+o(1)}
\]
with executed multiplicity and $h<1$.  Then depth-first
recursive execution decides every input in the covered typed domain in
\[
 T(N)\le N^{C+o(1)},\qquad
 C=\max\left\{1,r,\frac{v-h}{1-h}\right\},                \tag{2}
\]
\[
 S(N)\le N^{P+o(1)},\qquad P=\max\{p,h\}.                \tag{3}
\]
Same-scale internal alias nodes need neither a separate recursive controller
nor a shrinkage exponent.
\end{theorem}
\begin{proof}
Flatten the finite internal DAG into the implementation of $A_\tau$.  Its
work and retained state are already included in $r,p$; because it makes no
semantic call, it contributes no recursive term.  For $C\ge1$, the executed
boundary children satisfy
\[
 \sum_i |Q_i|^{C+o(1)}
 \le \left(\max_i|Q_i|\right)^{C-1+o(1)}
      \sum_i|Q_i|
 \le N^{v-h+hC+o(1)}.
\]
Inequality $C\ge v-h+hC$ is equivalent to the third term of (2).  Local phase work
gives $C\ge r$, and $C\ge1$ absorbs linear conventions and finite levels.
Equation (1) licenses induction on strictly decreasing boundary length.

Execute children one at a time.  The materialized child costs
$N^{h+o(1)}$ space; phase state costs $N^{p+o(1)}$.  Descendant scales form a
geometric exponent chain, so retained ancestor frames and child workspaces
sum to $N^{\max\{p,h\}+o(1)}$.  No same-size virtual-source buffer is retained.
\end{proof}

\paragraph{Typed form.}
With type-dependent bounds, a sufficient joint system is
\[
 C_\tau\ge
 \max\!\left\{1,r_\tau,
 \max_{e:\tau\to\sigma}
   \bigl(v_e-h_e+h_eC_\sigma\bigr)\right\},
\]
\[
 P_\tau\ge
 \max\!\left\{p_\tau,
 \max_{e:\tau\to\sigma}\{h_e,h_eP_\sigma\}\right\}.
\]
These inequalities are imposed only on completed boundary edges.  A
same-scale internal edge is legal only after its source-only computation has
been flattened into $r_\tau,p_\tau$ and proved well founded.

\section{What this recovers for the double wrapper}
Construction8 gives an actual internal chain
\[
 G(F)\longrightarrow \text{read-only alias of }F
 \longrightarrow \textsc{PIVOT-FACTOR1}
 \longrightarrow (Q_i)_i.
\]
The first arrow is same scale, acyclic, source-checkable and contains no SAT
call.  Therefore $F$ itself need not be a recursively hereditary guarded
instance.  Theorem~\ref{thm:alias} applies with
$\alpha=1,a=1,b=g=s=0$.  The finite measured ratios below one do not change
this asymptotic fact: $|F|=\Theta(|G(F)|)$.

Use the ordered portfolio PIVOT followed by iterative pure-literal DECIDE.
PIVOT keeps
\[
 (r,v,h,p,g)=(6/5,11/10,1/2,1/2,0)
\]
and the zero-child pure decision is linear with polylogarithmic state.
The union therefore keeps $(R,V,H,B)=(6/5,11/10,1/2,1/2)$.

\paragraph{Correction to Construction8.}
For its lifted-cross source $T_s$, PIVOT returns UNKNOWN and the model set is
non-affine, exactly as proved there.  But every $a_i,b_j$ is pure positive.
After guard stripping, setting those variables true deletes every source
clause.  Thus the two-member portfolio DECIDEs SAT with zero children.
$G(T_s)$ is a counterexample to PIVOT-plus-affine coverage, not to the
alias-lifted current portfolio as a whole.  This correction narrows rather
than closes the open set.

A sign-balanced proper lift with both $(x\vee a_i)$ and
$(x\vee\neg a_i)$ branches and both negative-$x$ analogues defeats both
members in the bounded controls.  More generally, no collective all-input
coverage theorem is proved.  If every emitted component $Q_i$ were covered
at the next ordinary boundary, Theorem~\ref{thm:closure} would give
$(C,P)=(6/5,1/2)$ and preserve conditional elimination of $K$.  Construction6's
flat marked parent instead emits an unmarked child.  These are uncovered
completed boundaries, not alias failures.  A scale-one internal cycle is a
different failure: coverage cannot repair the lack of finite flattening.

\section{Exact recovery audit}
\begin{center}
\begin{tabular}{@{}lll@{}}
\toprule
Record & Reusable exact result & Does not provide\\
\midrule
RECURCLOSE1 & strict recursive shrinkage closure & phase construction\\
CYCLEGAIN1 & same-scale edges inside a well-founded phase & a finite interface\\
STREAMCOMP1 & charge for implicit child access & an accessor implementation\\
ORBITBATCH1 & deduplication of declared field addresses & call/output deduplication\\
COMBINED-REPLAY1 & exact terminating all-input semantics & favorable resources\\
Recovery7 & aggregate-volume fixed point & boundary coverage\\
Construction8 & actual alias plus favorable partial boundary & all-input descendants\\
\bottomrule
\end{tabular}
\end{center}

Thus several previously open accounting items are solved, but their domains
do not silently unify.  COMBINED-REPLAY1 ranges over every finite CNF but may
have exponential Shannon multiplicity and linear state.  Construction8 has
the favorable ledger but only on guarded sources whose decoded source passes
PIVOT.  Theorem~\ref{thm:closure} composes these facts only when the same
implementation satisfies the boundary premise; it cannot move a conclusion
between their different domains.

\section{Literature/model check}
Classical SAT self-reducibility is 2-disjunctive and length decreasing:
$F$ is satisfiable iff one of $F[x=0],F[x=1]$ is satisfiable.  This is an
ordinary recursive boundary with two semantic children, not a source-only
alias, and both executions must be charged on UNSAT
(\url{https://arxiv.org/abs/1902.08299}).  The distinction is therefore
consistent with, rather than stronger than, the standard notion.

The downstream lower-bound source checked for this run is the
Fortnow--Lipton--van Melkebeek--Viglas deterministic random-access-machine
SAT theorem
(\url{https://pages.cs.wisc.edu/~dieter/Research/sat-ts.html}).
This report does not promote (2)--(3) to a separation and does not substitute
the retired unrestricted cosine inequality for the exact primary theorem.

\section{Bounded validation and remaining target}
The deterministic companion test uses a 24-second alarm.  It checked 144
exact rational fixed points, three saturated-access controls, all 64 directed
graphs on three internal types (25 DAGs and 39 rejected cycles), four
boundary-coverage controls, six PIVOT winners, eight repaired pure-literal
winners, and seven all-fail sign-balanced controls with zero child calls in
$3.436593$ seconds.  Digest:
\texttt{737ab8b3634f5409a22bd8a3edda1f7f930a4538d5084b97649afe9d85eada8a}.
These checks are \textsc{Computational}; they do not prove the theorem.

The open premise is now narrower:
\[
\boxed{\text{Every emitted ordinary-SAT boundary child must re-enter a
favorable controller, with common charged failure bounds.}}
\]
The next target is a successor-closed controller for the sign-balanced
all-fail family, or a proof that a finite portfolio collectively covers all
emitted component children.  P versus NP remains \textsc{Open}.

% END_OF_VERIFIED_REPORT::RECOVERY-STRENGTHEN::RECOVERY-STRENGTHEN8
\end{document}
